\begin{recipe}
[% 
    preparationtime = {\unit[15]{min}},
    portion = {\portion{1}},
    %source = {},
    calory = {950kcal | 35g Protein | 75g KH | 55g Fett},
    price = {4,80€},
    created = {17.03.2026},
    %rating = {1},
]
{Burger}

\graph{% pictures
    big=pic/nobrainer/01_burger
}

\introduction{%
Man braucht eigentlich kein Rezept für Burger. Aber falls doch: hier ist eins.
}

\ingredients{
    2 & Burger Buns\index{Getreide, Pasta \& Brot!Burger Buns}\\
    \nicefrac{1}{2} Packung & Vegane Burger-Trockenmischung\index{Tofu \& Fleischalternativen!Burger-Patty}\\
    2 Scheiben & Käse\index{Milchprodukte \& Alternativen!Käse}\\
    \nicefrac{1}{2}--1 & Avocado\index{Obst!Avocado}\\
    1 & Tomate\index{Gemüse \& Pilze!Tomaten}\\
    & Gurken (Essiggurken)\index{Vorrat \& Konserven!Essiggurken}\\
    & Salat\index{Gemüse \& Pilze!Salat}\\
    & Burgersauce\index{Öle, Saucen \& Würzmittel!Burgersauce}\\
    & Ketchup\index{Öle, Saucen \& Würzmittel!Ketchup}
}

\preparation{%
    \step%
    Burger Patties nach Packungsanweisung zubereiten.
    
    \step%
    Buns optional kurz anrösten.
    
    \step%
    Avocado in Scheiben schneiden, Tomate schneiden.
    
    \step%
    Burger nach Belieben belegen: Sauce, Salat, Patty, Käse, Gemüse, nochmal Sauce.
    
    \step%
    Zusammenklappen und essen.
}

\suggestion{%
    Wenn du hier ein Rezept brauchst, hast du ganz andere Probleme.
}

\end{recipe}